% =====================================================================
%  28.1  Rooted tree isomorphism and the AHU canonical labelling.
%  Exports: sec:c22-ahu  def:c22-iso  def:c22-canlab  thm:c22-ahu
%  lst:c22-ahu  fig:c22-twin  obs (2, inline)
%  WORKED section only -- no solvedproblem (the chapter's labs are 28.2
%  CF 601D and 28.3 CF 1800G).  DISCHARGES the Ch.27 (\chref{c21}) boundary
%  hand-off "which subtree EQUALS which".  Folds SHAPES with Ch.13's
%  bottom-up recursion; the string-encoding blow-up and the arbitrary-root
%  trap are owned HERE and only referenced later.
% =====================================================================
\section{Rooted Isomorphism: Sorting Shapes Bottom-Up}\label{sec:c22-ahu}

\begin{previouslybox}[Ch.~27 ended by asking the question this section answers]
Chapter~27's method boundary (\(\secref{c21-boundary}\)) closed with a parting
case Mo cannot compress: batches of queries whose answer depends on
\emph{which subtree equals which} -- equality of \emph{structure}, not of any
arithmetic statistic over values.  A frequency table has no entry for "shape";
distinct-count, \(\sum\mathrm{freq}^2\), nothing telescopes toward it.  The
previous chapter folded \emph{numbers} bottom-up (Ch.~13, \chref{c10}); this
chapter keeps the recursion and changes the folded object to the
\textbf{shape itself}.  One bottom-up pass will give every subtree a name such
that \emph{names are equal exactly when shapes are equal} -- the move the
boundary promised.
\end{previouslybox}

\begin{topicintro}
Two rooted trees are \term{isomorphic} when a bijection of vertices maps the
root to the root and preserves the parent--child relation: labels are noise,
child \emph{order} is noise, only the unlabelled unordered structure remains.
The AHU algorithm (Aho, Hopcroft, Ullman -- the name it carries in the
literature this chapter borrows from) decides the question in linear-ish time
with one idea: \textbf{label leaves first, then label each node by the sorted
tuple of its children's labels, giving each freshly-seen sorted tuple a fresh
integer}.  Identical shapes are forced to identical labels; distinct shapes get
distinct labels because the tuples themselves already differ one level down.
Root labels equal \(\iff\) isomorphic, and the same table of ids classifies
\emph{all} subtrees of one tree at once (\(\thmref{c22-ahu}\)).

\textbf{Prerequisites.} The bottom-up fold of Ch.~13 (\chref{c10}) -- same
recursion order, new payload.  Nothing else; the centre machinery that lifts
this to \emph{unrooted} trees is Ch.~12's and arrives in \(\secref{c22-centre}\).

\textbf{After this section you can}
\begin{itemize}[nosep]
  \item state the isomorphism definition for rooted trees and spot which statements are secretly asking for it; \hfill\lv{3}
  \item run the canonical-label fold by hand on a 9-node tree and read off every subtree's shape name; \hfill\lv{3}
  \item code the interned-tuple labeller with a \texttt{map<vector<int>,int>}; \hfill\lv{4}
  \item explain why encoding shapes as \emph{strings} blows up while encoding them as \emph{interned ids} does not. \hfill\lv{4}
\end{itemize}
\end{topicintro}

% ---------------------------------------------------------------
\subsection{Canonical labels: sort, intern, compare}
% ---------------------------------------------------------------

\begin{definition}{Rooted isomorphism}{c22-iso}
Rooted trees \((T,r)\) and \((T',r')\) are \term{isomorphic}, written
\((T,r)\cong(T',r')\), if a bijection \(f:V(T)\to V(T')\) exists with
\(f(r)=r'\) and \((u,v)\) an edge of \(T\) iff \((f(u),f(v))\) an edge of \(T'\).
Equivalently: \(f(r)=r'\) and \(f\) maps the children of each node bijectively
onto the children of its image.  Child order is \emph{not} part of the structure.
\end{definition}

\begin{definition}{Canonical label \(\mathrm{cl}(u)\) by bottom-up interning}{c22-canlab}
Process nodes in postorder.  For a leaf, \(\mathrm{cl}(u)\) is the id interned for the
empty tuple.  For an internal node, sort the multiset
\(\{\mathrm{cl}(c):c\in\children(u)\}\) into a tuple \(\tau_u\) and set \(\mathrm{cl}(u)\) to the
id of \(\tau_u\): a fresh counter \(next\) hands out
\(\mathit{next}++\) the first time a tuple is seen, and the map
\(\mathit{tid}[\tau]\) recalls the old id thereafter.  Two tuples get the same id
\emph{by construction} iff they are equal as sorted sequences.
\end{definition}

\begin{theorem}{AHU: \((T,r)\cong(T',r')\) iff the interned root labels match}{c22-ahu}
One bottom-up pass with a shared id table labels both trees.  If the roots share
a label the trees are isomorphic (unfold the tuple equality into a bijection
child-by-child); if they are isomorphic the roots share a label (induction on
height: isomorphic nodes have isomorphic child multisets, hence equal sorted
tuples, hence equal ids).
\end{theorem}
\begin{proofbox}
\emph{(<=)} Let the root tuples of \(T\) and \(T'\) be the same sorted sequence
\((\ell_1\le\cdots\le\ell_k)\).  Children may \emph{repeat} a label (two
identical subtrees), so pair by multiplicity: the tuples are equal as multisets,
hence there is a bijection between \(\children(r)\) and \(\children(r')\) pairing
equal labels.  By induction each paired pair of children roots isomorphic
subtrees (equal id \emph{means} isomorphic -- the induction hypothesis applied to
the interned tuple); composing the child-level bijections under the root map
gives the claimed \(f\). \emph{(=>)} Induct on \(\height\).  Leaves: both
tuples are empty, both get the empty tuple's id.  For height \(h\): an isomorphism
\(f\) maps \(\children(u)\) bijectively onto \(\children(f(u))\) pairing
isomorphic children; by induction paired children carry equal labels; so the two
sorted tuples are identical sequences and intern to one id. \(\Box\)
\end{proofbox}

\begin{lstlisting}[caption={The interned-tuple labeller: one postorder DFS fills cl[1..n] for a rooted tree; comparing two trees (or two subtrees) is one integer compare. ASCII only, identifiers without underscores}, label={lst:c22-ahu}]
#include <bits/stdc++.h>
using namespace std;
const int MAXN=200005;
int n, nextId=1;
vector<int> adj[MAXN], kids[MAXN];
map<vector<int>, int> tid;          // sorted tuple -> interned id
int cl[MAXN];                       // canonical label of each node
int parentArr[MAXN];
void rootify(int u, int p){
    parentArr[u]=p;
    for(int v: adj[u]) if(v!=p){ kids[u].push_back(v); rootify(v,u); }
}
int label(int u){
    vector<int> t;
    for(int v: kids[u]) t.push_back(label(v));
    sort(t.begin(), t.end());
    map<vector<int>,int>::iterator it = tid.find(t);
    if (it != tid.end()) return cl[u]=it->second;
    int id = nextId++;
    tid[t]=id;
    return cl[u]=id;
}
// usage: rootify(1,0); label(1);  then (T,1) and (T',1') iso iff cl==cl'
//        (with ONE shared tid/nextId across both trees!)
\end{lstlisting}

\begin{insightbox}[Interned ids, not strings -- the \(O(n\log n)\) vs \(O(n^2)\) cliff]
The textbook version encodes a tree as a balanced-parentheses string
\((\,\text{child strings sorted}\,)\).  Correct, but a string \emph{copy} at every
ancestor costs \(\Theta(\size)\) per node: \(\bigO{n^2}\) bits.  Interning replaces
each string by a small integer \emph{names} for the same object: the tuple of
child names has total length \(n-1\) summed over all nodes, so sorting all tuples
is \(\bigO{n\log n}\) and the map pays \(O(\text{tuple length}\cdot\log\#\text{tuples})\)
per distinct tuple -- still near-linear.  Names-vs-strings is the whole
optimization, and it is exactly the trick hashing will re-use from the other side.
\end{insightbox}

\begin{tracebox}[The AHU fold on the running tree \(R\) (opener, figure \(\figref{c22-tree}\))]
Postorder (children in listing order): \(5,6,2,7,8,3,9,4,1\).
Leaves \(5,6,7,8,9\): the empty tuple interns first \(\Rightarrow\) \(\mathrm{cl}=\mathbf{1}\).
Node \(2\): children \((5,6)\to\) tuple \((1,1)\) interns second \(\Rightarrow\mathrm{cl}=\mathbf{2}\).
Node \(3\): children \((7,8)\to\) tuple \((1,1)\) again \(\Rightarrow\mathrm{cl}=2\) (recalled).
Node \(4\): child \((9)\to\) tuple \((1)\Rightarrow\mathrm{cl}=\mathbf{3}\).
Root \(1\): children \((2,3,4)\to\) labels \((2,2,3)\Rightarrow\mathrm{cl}=\mathbf{4}\).
Readings: subtree\((2)\cong\) subtree\((3)\) (both \(\mathrm{cl}=2\)); subtree\((4)\not\cong\)
subtree\((2)\) (\(3\ne2\)); the whole tree has name \(4\).  Four distinct shapes on
nine nodes: the leaf, the 2-leaf star, the 1-leaf chain, and the root's
star+star+chain. \(\Box\)
\end{tracebox}

\subsection{Comparing two trees -- and the shared-table trap}
% ---------------------------------------------------------------

Two separate trees are compared by labelling \emph{both against one shared}
\(\mathit{tid}\) table and reading the roots' ids.  Two subtrees of one tree are
compared by a single labelling and one integer compare, \(O(1)\) per query after
the \(O(n)\) pass -- Ch.~27's ``which subtree equals which'', answered.  For the
\emph{unrooted} question the given roots are not part of the input, and equal
rooted ids at arbitrary vertices prove nothing; \(\secref{c22-centre}\) fixes
that with the canonical root.  The twin \(R'\) below previews both verdicts.

\begin{figure}[H]
\centering
\begin{tikzpicture}[diagram, scale=0.95, transform shape,
   tnode/.style={circle, draw=primaryblue!70, fill=primaryblue!8, line width=0.7pt,
                 minimum size=6.6mm, font=\scriptsize\bfseries, inner sep=1pt},
   sub/.style={circle, draw=successgreen!70, fill=successgreen!14, line width=1.0pt,
             minimum size=6.6mm, font=\scriptsize\bfseries, inner sep=1pt}]
  \node[sub]   (m1) at (0,-4.3) {1};
  \node[tnode] (m2) at (0,-1.5) {2};
  \node[tnode] (m3) at (-2.6,-3.0) {3};
  \node[tnode] (m4) at (2.6,-3.0) {4};
  \node[tnode] (m5) at (0,0) {5};
  \node[sub]   (m6) at (-1.4,-4.3) {6};
  \node[tnode] (m7) at (-4.0,-4.0) {7};
  \node[tnode] (m8) at (-3.4,-1.2) {8};
  \node[tnode] (m9) at (2.6,-4.3) {9};
  \draw (m5)--(m2) (m5)--(m6) (m2)--(m3) (m2)--(m4) (m3)--(m7) (m3)--(m8) (m4)--(m9) (m1)--(m5);
  \node[annot, font=\scriptsize, anchor=north west, text width=9.8cm] at (-4.6,1.2)
    {\(R'\): edges 1-5, 2-5, 2-3, 2-4, 3-7, 3-8, 4-9, 5-6 -- built as \(\sigma(R)\), \(\sigma=(1\,2\,5)\)};
\end{tikzpicture}
\caption{\(R'\), the unrooted twin of \(R\), constructed \emph{by definition}:
apply the renaming \(\sigma=(1\;2\;5)\) (cycle \(1\mapsto2,\,2\mapsto5,\,5\mapsto1\), all
other vertices fixed) to every edge of \(R\):
\(1\!-\!2\mapsto2\!-\!5\)✓, \(1\!-\!3\mapsto2\!-\!3\)✓, \(1\!-\!4\mapsto2\!-\!4\)✓,
\(2\!-\!5\mapsto5\!-\!1\)✓, \(2\!-\!6\mapsto5\!-\!6\)✓, \(3\!-\!7\mapsto3\!-\!7\)✓,
\(3\!-\!8\mapsto3\!-\!8\)✓, \(4\!-\!9\mapsto4\!-\!9\)✓ -- exactly the listed edge set.
A twin built as \(\sigma(R)\) needs no re-derivation: \(\sigma\) \emph{is} the
witnessed isomorphism.  Yet \emph{rooted} labels at vertex 1 disagree:
\(R\) at root 1 has name \(4\); labelling \(R'\) at root 1 (rooting gives
\(1\!\to\!5\!\to\!\{2,6\}\), \(2\!\to\!\{3,4\}\), \(3\!\to\!\{7,8\}\),
\(4\!\to\!9\); postorder \(6,7,8,3,9,4,2,5,1\), on \(R\)'s table):
leaves \(6,7,8,9\to\mathrm{cl}=1\); \(3\to(1,1)\to\mathrm{cl}=2\) and
\(4\to(1)\to\mathrm{cl}=3\), both \textbf{recalled}; then
\(2\to(2,3)\to\) \textbf{fresh} \(\mathrm{cl}=\mathbf{5}\);
\(5\to(1,5)\to\) fresh \(\mathrm{cl}=\mathbf{6}\); root
\(1\to(6)\to\) fresh \(\mathrm{cl}=\mathbf{7}\).
\(R'\)'s root wears \(7\ne4\).  Same unrooted tree, different rooted shape --
because \(R'\)'s vertex 1 is \(R\)'s vertex \(\sigma^{-1}(1)=5\), a leaf, now
wearing the root's hat.  The verdicts 28.3 will print: rooted-at-given-1
\textbf{NO}, centre-rooted (\(R\)'s centre 1, \(R'\)'s centre 2) \textbf{YES}.}
\label{fig:c22-twin}
\end{figure}

\begin{warningbox}[Three ways the fold goes wrong]
\begin{enumerate}[nosep,leftmargin=1.6em]
  \item \textbf{Separate id tables per tree.} Interning tree A then tree B with
        fresh counters makes B's ``leaf id'' different from A's ``leaf id'' and
        every compare lies.  One shared \(\mathit{tid}\) (or hash everything,
        \(\secref{c22-hash}\)).  Symptom: identical trees return NO.
  \item \textbf{Comparing rooted ids to answer an unrooted question.} \(R\) at
        root 1 and \(R'\) at root 1 have different names -- yet the unrooted trees
        are identical.  Rooted ids decide rooted questions only until 28.3
        canonicalizes the root.
  \item \textbf{Sorting node \emph{indices} instead of child \emph{labels}.} The
        tuple must be sorted by \(\mathrm{cl}\) values; sorting adjacency lists by vertex
        number makes the tuple depend on naming -- reintroducing by the back door
        the labels you were trying to erase.
\end{enumerate}
\end{warningbox}

\begin{hardeningbox}
\begin{itemize}[nosep,leftmargin=1.4em]
  \item \textbf{Iterate by height, not by recursion.} For \(n=2\cdot10^5\) the
        postorder DFS may stack-overflow; process nodes in decreasing height
        (bucket by height from the leaf peeling of Ch.~12) -- same fold, no stack.
  \item \textbf{\(O(n)\) is reachable.} Radix-sort each child-label multiset
        (labels are in \([0,n]\)) and bucket tuples by length, then lexicographically
        -- the linear-time AHU.  In contests \texttt{map<vector<int>,int>} at
        \(\bigO{n\log n}\) amortized by tuple lengths passes comfortably; swap in
        \texttt{unordered\_map} with a vector hash only if the constant bites.
  \item \textbf{Re-interning across test cases.} With \(t\le10^4\) cases, clearing
        \(\mathit{tid}\) per case is fine \emph{if} both trees of one case share
        the case's table (pitfall 1 again); carrying one table across cases also
        works and costs nothing.
\end{itemize}
\end{hardeningbox}

\begin{recallbox}[28.1 in one breath]
\emph{Shape names by sort-and-intern: process postorder, hand each node the id of
its \textbf{sorted} child-label tuple, share the table between everything you
compare -- then equality of shapes is equality of integers
(\(\thmref{c22-ahu}\), \(\coderef{c22-ahu}\)).  On \(R\): leaf \(=_{\mathrm{cl}}1\),
star \(=_{\mathrm{cl}}2\), chain-of-two \(=_{\mathrm{cl}}3\), root \(=_{\mathrm{cl}}4\);
the twin \(R'\) rooted at 1 wears name \(7\) while \(R\) rooted at 1 wears 4 --
same unrooted tree, different rooted tree, the gap 28.3 closes.  Interning beats parentheses strings the way
pointers beat photocopies (\(\bigO{n\log n}\) vs.\ \(\bigO{n^2}\)).}
\selftest{
  \item Why must the child labels be sorted before interning? \emph{(Child order is not structure; the sorted tuple is the order-free name -- unsorted, two isomorphic trees could intern different tuples.)}
  \item One tree, every subtree pair queried: total cost? \emph{(One labelling O(n log n), then O(1) per compare.)}
  \item \(R\) vs \(R'\): name at \(R\)'s root is 4, at \(R'\)'s root 7 -- does that disprove isomorphism? \emph{(No -- the twins are the same unrooted tree; only the ROOTED verdict differs; unrooted iso needs centre-rooted comparison, 28.3.)}
}
\end{recallbox}
